\subsection{局部紧拓扑向量空间}

定理 3. --- 设 K 是完备非离散赋值除环。若 E 是 K 上的一个 Hausdorff 拓扑向量空间，且 E 中 0 的某个邻域 V 是预紧的，则 E 是有限维的。若 \(E \neq \{0\}\)，则 K 与 E 都是局部紧的。

在证明第一个论断时，只需考虑 E 是完备的情形；因为 E 是其完备化 \(\hat{E}\) 的一个处处稠密的子空间，且 V 的闭包 \(\overline{V}\)（取在 \(\hat{E}\) 中）是紧的，并且是 \(\hat{E}\) 中 0 的一个邻域（GT, III, § 3.4, prop. 7）。

于是可以假设 E 中 0 有一个紧邻域 V。设 \(\alpha \in \mathbf{K}\) 满足 \(0 < |\alpha| < 1\)；则存在有限多个点 \(a_i \in \mathbf{V}\) 使得

\[
\mathrm{V} \subset \bigcup_ {i} (a _ {i} + \alpha \mathrm{V}).
\]

设 M 是 E 中由诸 \(a_{i}\) 生成的有限维子空间；它在 E 中是闭的（I, 第 14 页, cor. 1）。在 Hausdorff 拓扑向量空间 E/M 中，V 的典范像 W 是 0 的一个紧邻域，满足 \(W \subset \alpha W\)；由此 \(\alpha^{-1} W \subset W\)，并且对 n 作归纳可得：对每个正整数 n 有 \(\alpha^{-n} W \subset W\)。由于 W 是吸收的，我们断言 W = E/M；从而 E/M 是紧的。因此，要完成定理中第一个论断的证明，只需证明下述引理。

引理 1. --- 非离散赋值除环上的每个紧拓扑向量空间 E 都恰好是集合 \(\{0\}\)。

由于 E 是完备的，可以假设 K 是完备的（I, 第 6 页）。若 \(E \neq \{0\}\)，则 E 含有一条在 E 中闭的直线（I, 第 14 页, cor. 1），从而这条直线是紧的。这条直线同构于 \(K_{s}\)（I, 第 12 页, prop. 2），因而 K 必是紧的。但映射 \(\xi \mapsto |\xi|\)（K 到 R 中的映射）是连续的，故 K 的像必有界；另一方面，存在 \(\gamma \in K\) 使 \(|\gamma| > 1\)，而集合 \(|\gamma^{n}| = |\gamma|^{n}\)（\(n \in N\)）是无界的。这一矛盾表明 \(E = \{0\}\)。

为证明定理中的第二个论断，若 \(E \neq \{0\}\)，则由定理的第一部分，E 同构于 \(K_{s}^{n}\)，其中 n \textgreater{} 0；而 K 是完备的，故 E 也是完备的，从而 E 是局部紧的。但 \(K_{s}\) 同构于 E 中的一条直线（I, 第 12 页, prop. 2），后者在 E 中必然是闭的（I, 第 14 页, cor. 1）；由此推出 K 是局部紧的。

注记. --- 当 K 是离散除环时，th. 3 的结果不再成立，正如把 R（取通常拓扑）视为离散域 Q 上的拓扑向量空间这一例子所示。

